\documentclass[article]{jss}

% Working skeleton: replace every TODO before submission.
% Leave publication metadata (volume, issue, DOI, dates) to the editors.
\newcommand{\todo}[1]{\textbf{TODO:} #1}

\author{Brandon Greenwell\\Affiliation to be confirmed}
\Plainauthor{Brandon Greenwell}
\title{\pkg{guide-rs}: GUIDE-Style Classification and Regression Trees
  in \proglang{Rust}}
\Plaintitle{guide-rs: GUIDE-Style Classification and Regression Trees in Rust}
\Shorttitle{GUIDE-Style Trees in \proglang{Rust}}
\Abstract{
  We present \pkg{guide-rs}, an open-source implementation of GUIDE-style
  classification and regression trees in \proglang{Rust}. The software
  provides a reusable library and a command-line interface for fitting,
  inspecting, saving, and applying models. Its pure \proglang{Rust} core
  supports native execution on Linux, macOS, and Windows without requiring
  the original GUIDE executable or external numerical libraries. A
  complementary \proglang{Python} wrapper provides
  \pkg{scikit-learn}-compatible estimators. The paper focuses on
  association-based variable selection, constant classification and
  regression trees, multiple-linear regression leaves, and pruning.
  The contribution centers on source availability, reusable interfaces,
  and integration into reproducible statistical workflows.
  \todo{Add the completed application's findings and scoped validation
  results. This draft reports no empirical results.}
}
\Keywords{recursive partitioning, variable selection, model-based trees,
  reproducible research, \proglang{Rust}, \proglang{Python}}
\Plainkeywords{recursive partitioning, variable selection, model-based trees,
  reproducible research, Rust, Python}
\Address{
  Brandon Greenwell\\
  Affiliation and postal address to be confirmed\\
  E-mail: \email{greenwell.brandon@gmail.com}\\
  URL: \url{https://github.com/bgreenwell/guide-rs}
}

\begin{document}

\section{Introduction}\label{sec:introduction}

Classification and regression trees express prediction rules through
successive partitions of the predictor space. Their statistical behavior
depends on the procedures used to select variables, choose splits, fit
terminal-node models, and prune the resulting tree. GUIDE (generalized,
unbiased, interaction detection and estimation) separates variable selection
from split-point search and incorporates tests for local interactions
\citep{Loh:2002,Loh:2009}. These methods provide the statistical foundation
for \pkg{guide-rs}.

Access to a statistical method also depends on how its implementation is
distributed and used. GUIDE has historically been available free of charge
as closed-source software \citep{Rusch+Zeileis:2014}. Its current distribution
provides executables for Linux, macOS, and Windows and supports both
interactive and batch operation \citep{GUIDE:software}. An openly available
implementation additionally enables researchers to inspect numerical
decisions, rebuild the software, adapt its methods, and incorporate its
components into other applications.

We present \pkg{guide-rs}, an independent, MIT-licensed implementation of
a subset of GUIDE-style methods. Its central contribution is to make these
methods available as open, reusable statistical computing infrastructure.
The intended users include analysts conducting scripted studies, developers
embedding tree models in applications, and researchers examining or extending
the implementation. The original GUIDE publications establish the methods;
this article describes their implementation, interfaces, and evaluation.

\subsection{Software contribution}\label{sec:contribution}

The contribution has four connected elements. First, the source code makes
the implementation available for inspection, modification, and redistribution.
Statistical routines and implementation choices can be examined alongside
tests and documented numerical conventions.

Second, a reusable \proglang{Rust} library and an explicit command-line
interface (CLI) support both embedded and scripted use. Library users fit
models and obtain predictions directly in memory. CLI users supply data and
structured configuration to separate commands for fitting, inspection, and
prediction, with fitted models saved in JSON format. These interfaces allow
applications to call the implementation directly and scripts to specify
their operations without an interactive dialogue. A worked example will
illustrate the complete workflow and its inputs and outputs.

Third, the pure \proglang{Rust} computational core supports native deployment
and is tested on Linux, macOS, and Windows. The CLI requires neither the
original GUIDE executable nor a \proglang{Python} interpreter, and the core
has no external BLAS or LAPACK dependency. The source distribution permits
users to build the library and executable for supported targets. This
combination of source availability and reusable interfaces complements the
platform-specific executables already available for GUIDE.

Fourth, a complementary \proglang{Python} wrapper exposes the same core
through estimators compatible with \pkg{scikit-learn}
\citep{Pedregosa+etal:2011}. Familiar fitting and prediction methods,
pipelines, and model-selection tools provide a convenient entry point for
existing \proglang{Python} users. Statistical fitting and configuration
validation remain in the shared \proglang{Rust} implementation.

The implementation is accompanied by frozen comparisons with native GUIDE,
independent numerical checks, and documented behavioral differences.
These provide evidence for specific capabilities and numerical objectives.
Section~\ref{sec:validation} will distinguish reference agreement from
independently checked correctness and state the limits of each comparison.

\subsection{Scope}\label{sec:scope}

The methodological treatment and evaluation focus on association-based
variable selection and split search, constant classification and
least-squares regression trees, multiple-linear least-squares leaves, and
pruning. Numeric and categorical predictors and missing-value handling are
considered within these model families. The software presentation leads
with the \proglang{Rust} library and CLI and then introduces the
\proglang{Python} wrapper.

Other implemented capabilities include additional losses and node models,
bagging and forests, and variable importance. Their coverage and validation
status will be summarized separately. Gradient boosting is an experimental
\pkg{guide-rs} extension. Survival, treatment-effect, longitudinal, and
multiresponse methods are outside the implemented scope described here.
\pkg{guide-rs} remains experimental: it implements a subset of GUIDE's
capabilities and has known numerical and behavioral differences from native
GUIDE. Feature availability, reference agreement, and numerical correctness
are separate aspects of the evaluation.

\todo{Complete the related-software comparison, including conventional CART
implementations and the conditional inference and model-based trees in
\pkg{partykit} \citep{Hothorn+Zeileis:2015}. Pin the release and match each
final claim to its validation evidence.}

\section{Statistical methods}\label{sec:methods}

\todo{Specify the statistical procedure and its assumptions.}

\subsection{Variable selection and split search}
\todo{Describe association-based variable selection, interaction tests,
and subsequent split-point search. Specify the implemented tests, stopping
rules, and treatment of numeric, categorical, and missing predictors.
Distinguish published methods from implementation choices.}

\subsection{Node models and pruning}
\todo{Specify the response families and node models covered by the paper,
their fitting objectives, and pruning procedures. Begin with constant
classification and regression and multiple-linear least-squares leaves;
expand the scope only when the corresponding evidence is ready.}

\section{Software design and interfaces}\label{sec:software}

\todo{Explain the flow from data and configuration through fitting to
prediction. Describe the pure \proglang{Rust} core, shared validation,
saved models, deterministic seeds, and parallel computation. Explain
design decisions in terms of their benefit to statistical users.}

\subsection[Rust library]{\proglang{Rust} library}
\todo{Demonstrate the primary library API: data construction, configuration,
fitting, prediction, and model persistence. Explain application integration
and validation errors; leave API inventories to the software documentation.}

\subsection{Command-line workflow}
\todo{Demonstrate a scripted CSV/JSON workflow: fit, inspect, save, and predict.
Explain code in prose; use \code{CodeInput} and \code{CodeOutput} for examples.}

\subsection[Python convenience wrapper]{\proglang{Python} convenience wrapper}
\todo{Demonstrate convenient fitting, prediction, and tuning with
\code{GuideRegressor} or \code{GuideClassifier}. Explain how the
scikit-learn-compatible wrapper delegates to the core and shares validation.}

\section{Validation and software comparisons}\label{sec:validation}

\todo{Summarize frozen native comparisons, independent numerical checks,
and statistical selection experiments for the pinned release. Separate
agreement with native GUIDE from independently established correctness.
Report known differences and the precise scope of each conclusion.}

\todo{Compare relevant capabilities and workflows with native GUIDE,
conventional CART implementations, and conditional inference and model-based
trees such as \pkg{partykit} \citep{Hothorn+Zeileis:2015}.
Prespecify any empirical comparisons, tuning budgets, metrics, and runtime
conditions. Report uncertainty and limitations; avoid an extensive
simulation study or unsupported claims of universal unbiasedness.}

\section{An applied example}\label{sec:application}

\todo{Choose a non-trivial, redistributable real dataset and a substantive
question. Document provenance and preprocessing, then show the complete
analysis through the CLI, with interpretation and held-out evaluation.
Follow with a concise illustration of the equivalent estimator workflow
through the \proglang{Python} wrapper. Generate every reported
number, table, and figure from the replication workflow.}

\section{Discussion}\label{sec:discussion}

\todo{Summarize the demonstrated contribution, appropriate use cases,
limitations, and future work. Separate supported features from experimental
extensions. Do not imply support for unimplemented GUIDE families.}

\section*{Computational details}

\todo{Record the software release and commit, algorithm revision,
\proglang{Rust} and \proglang{Python} versions, dependencies, operating
system, hardware where relevant, seeds, and thread settings. Give the
installation and replication commands and an archival release identifier.
The repository is \url{https://github.com/bgreenwell/guide-rs}.}

\section*{Acknowledgments}

\todo{Confirm authorship, affiliations, funding, and acknowledgments.
Add any required disclosures.}

\clearpage
\bibliography{refs}

% Optional appendices go after the bibliography, beginning on a new page.
% Use descriptive section titles, not a section titled only "Appendix".

\end{document}
